\chapter{Conclusion}
\label{chap:conclusion}

This thesis examined whether uncertainty in an estimated MI standard error
can be represented through an MI-specific Welch--Satterthwaite correction
when comparing two independent categorical populations. The null is equality
of their population MI values, not independence within either population or
equality of the complete joint distributions.

The derivation starts from the pointwise-MI variance governing first-order
sampling error. A one-cell perturbation changes the joint probabilities,
the margins and the pointwise values. Differentiating the complete variance
functional supplies the sampling variability needed for component effective
degrees of freedom. The resulting Expanded Welch procedure retains Wald's
bias-corrected MI difference and plug-in standard error but evaluates the
shared statistic against a Student reference. This is a principled
moment-matching construction, not an exact finite-sample Student law.

The completed fixed-population experiment found both benefits and costs.
Expanded Welch can make a liberal Wald null rate closer to the nominal 5\%
level. Where Wald is already conservative, the heavier-tailed reference moves
rejection farther below that target. In sparse regimes, Expanded Welch also
returns fewer valid results; a low unconditional rate there is not reliable
evidence of good calibration. At ordinary positive MI and sufficiently large
samples, both methods approached the nominal null rate in the tested fixed
populations, while near-independence settings remained difficult even at
large finite sample sizes. The expanded call used about 1.75 times the
runtime of Wald in the measured implementation.

Near independence, the component degree of freedom is approximately \(nI\),
so Expanded Welch can impose a very heavy-tailed reference where the
first-order approximation remains weak. A supplementary diagnostic confirms
that errors can arise in both the estimated standard error and the numerator.
Changing only the reference distribution cannot correct both.

The evidence therefore supports Normal Wald as the general baseline. Expanded
Welch is only a conditional sensitivity analysis; future corrections must
address the numerator, standard error and population structure together rather
than changing the reference law alone.
